\documentclass[aspectratio=169]{beamer}

% ============================================================
% Preambulo: tema propio replicando el estilo visual de
% context/diseno_experimentos.pdf (banda de titulo lavanda,
% cajas de color redondeadas, numeracion "n / N") y la
% numeracion de secciones de context/MARIELALVAREZ_QNN.pdf
% (plan de trabajo).
% ============================================================

\usepackage[utf8]{inputenc}
\usepackage[T1]{fontenc}
\usepackage[spanish,es-tabla]{babel}
\usepackage{lmodern}
\renewcommand{\familydefault}{\sfdefault}
\usepackage{helvet}
\usepackage[table]{xcolor}
\usepackage{mdframed}
\usepackage{booktabs}
\usepackage{array}
\usepackage{tikz}
\usetikzlibrary{arrows.meta,positioning,shadows}
\usepackage{amssymb}

\usetheme{default}
\setbeamertemplate{navigation symbols}{}
\setbeamertemplate{itemize items}[circle]
\setbeamertemplate{enumerate items}[default]

% ---------- paleta (muestreada de diseno_experimentos.pdf) ----------
\definecolor{bmBand}{HTML}{D6D6F0}   % banda de titulo / caja de portada
\definecolor{bmIndigo}{HTML}{1E1E79} % caja "framework" / marco primario
\definecolor{bmBlue}{HTML}{1E50A0}   % caja "Optimizacion"
\definecolor{bmGreen}{HTML}{14823C}  % caja "Pregunta de investigacion"
\definecolor{bmAmber}{HTML}{B48200}  % caja "Principios" / "Control"
\definecolor{bmFillBlue}{HTML}{F2F2FF}
\definecolor{bmFillGreen}{HTML}{F2FFF2}
\definecolor{bmFillAmber}{HTML}{FFFEE6}
\definecolor{bmInk}{HTML}{1A1A1A}
\definecolor{bmGrey}{HTML}{6B6B6B}
\definecolor{bmRowAlt}{HTML}{ECECF7}

\setbeamercolor{normal text}{fg=bmInk}
\setbeamercolor{frametitle}{fg=bmInk,bg=bmBand}
\setbeamercolor{title}{fg=bmInk}
\setbeamercolor{structure}{fg=bmIndigo}
\setbeamerfont{frametitle}{series=\bfseries,size=\large}
\setbeamerfont{footline}{size=\tiny}
\setbeamerfont{normal text}{size=\small}

% ---------- banda de titulo (sin sombra, ancho completo) ----------
\setbeamertemplate{frametitle}{
  \vskip-1pt
  \begin{beamercolorbox}[wd=\paperwidth,sep=0.30cm,leftskip=0.35cm]{frametitle}
    \usebeamerfont{frametitle}\insertframetitle
  \end{beamercolorbox}
}

% ---------- pie de pagina: solo "n / N" en la esquina ----------
\setbeamertemplate{footline}{
  \hfill
  \raisebox{2pt}{\color{bmGrey}\scriptsize\insertframenumber\,/\,\inserttotalframenumber}
  \hspace{4pt}
  \vspace{4pt}
}

% ---------- cajas de color (replican las de la p.2 y p.6) ----------
\newcommand{\bmheading}[2]{\textbf{\color{#1}#2}\par\smallskip}

\mdfdefinestyle{bmIndigoStyle}{backgroundcolor=bmFillBlue,linecolor=bmIndigo,
  linewidth=0.8pt,roundcorner=6pt,innertopmargin=7pt,innerbottommargin=7pt,
  innerleftmargin=9pt,innerrightmargin=9pt,skipabove=6pt,skipbelow=6pt}
\mdfdefinestyle{bmBlueStyle}{backgroundcolor=bmFillBlue,linecolor=bmBlue,
  linewidth=0.8pt,roundcorner=6pt,innertopmargin=7pt,innerbottommargin=7pt,
  innerleftmargin=9pt,innerrightmargin=9pt,skipabove=6pt,skipbelow=6pt}
\mdfdefinestyle{bmGreenStyle}{backgroundcolor=bmFillGreen,linecolor=bmGreen,
  linewidth=0.8pt,roundcorner=6pt,innertopmargin=7pt,innerbottommargin=7pt,
  innerleftmargin=9pt,innerrightmargin=9pt,skipabove=6pt,skipbelow=6pt}
\mdfdefinestyle{bmAmberStyle}{backgroundcolor=bmFillAmber,linecolor=bmAmber,
  linewidth=0.8pt,roundcorner=6pt,innertopmargin=7pt,innerbottommargin=7pt,
  innerleftmargin=9pt,innerrightmargin=9pt,skipabove=6pt,skipbelow=6pt}

% cajas de cuatro colores predefinidos, listas para usar
\newmdenv[style=bmIndigoStyle]{indigobox}
\newmdenv[style=bmBlueStyle]{bluebox}
\newmdenv[style=bmGreenStyle]{greenbox}
\newmdenv[style=bmAmberStyle]{amberbox}

% ---------- tablas con banda alterna ----------
\newcommand{\bandedtable}{\rowcolors{2}{white}{bmRowAlt}}

% ---------- comando de encabezado de subseccion numerada (5.2 estilo) ----------
\newcommand{\subsec}[2]{\textbf{\color{bmIndigo}#1.} #2}

% ---------- resaltar el modelo ganador en tablas ----------
\newcommand{\winner}[1]{\textbf{#1}}

\title{Benchmark reproducible de QCNN bajo condiciones NISQ}
\subtitle{Presentacion de avance}
\author{Mariel Alvarez Salas}
\date{17 de septiembre de 2026}

\begin{document}

% ============================================================
% 1. PORTADA
% ============================================================
\begin{frame}[plain]
  \vspace{1.2cm}
  \begin{center}
  \begin{tikzpicture}
    \node[rounded corners=6pt, fill=bmBand, drop shadow,
          text width=0.86\textwidth, align=center, inner sep=14pt] {
      {\Large\bfseries Benchmark reproducible de QCNN bajo condiciones NISQ}\\[4pt]
      {\normalsize Presentacion de avance de investigacion}
    };
  \end{tikzpicture}
  \vspace{1.0cm}

  \begin{tabular}{r l}
    \textbf{Estudiante:} & Mariel Alvarez Salas \\
    \textbf{Director:} & Dr. Leonardo Ledesma Dominguez \\
    \textbf{Programa:} & Next Generation Scientists (NGS) \\
    \textbf{Fecha:} & 17 de septiembre de 2026 \\
  \end{tabular}
  \end{center}
\end{frame}

% ============================================================
% 2. AGENDA
% ============================================================
\begin{frame}{2. Agenda}
  \begin{columns}[T]
    \begin{column}{0.5\textwidth}
      \begin{enumerate}
        \setlength{\itemsep}{7pt}
        \item[3.] Contexto y pregunta de investigacion
        \item[4.] Diseno experimental
        \item[5.] Arquitecturas evaluadas
        \item[6.] Resultados E0A
        \item[7.] Quantum Advantage
      \end{enumerate}
    \end{column}
    \begin{column}{0.5\textwidth}
      \begin{enumerate}
        \setlength{\itemsep}{7pt}
        \item[8.] Validez de la aportacion (E0B)
        \item[9.] Resultados E1 (shots finitos)
        \item[10.] Proximos pasos
        \item[11.] Cierre
      \end{enumerate}
    \end{column}
  \end{columns}
\end{frame}

% ============================================================
% 3. CONTEXTO Y PREGUNTA DE INVESTIGACION
% ============================================================
\begin{frame}{3. Contexto y pregunta de investigacion}
  \begin{indigobox}
    \bmheading{bmIndigo}{Pregunta central}
    Que tan robustas son las conclusiones sobre el desempeno de una QCNN cuando se
    transita de una simulacion analitica ideal a condiciones NISQ realistas: numero
    finito de shots, ruido de hardware y conectividad limitada de qubits.
  \end{indigobox}

  \vspace{10pt}

  \begin{columns}[T]
    \begin{column}{0.5\textwidth}
      \begin{greenbox}
        \bmheading{bmGreen}{Lo que este proyecto no es}
        No busca proponer una arquitectura QCNN nueva. El punto de partida son cuatro
        arquitecturas ya publicadas.
      \end{greenbox}
    \end{column}
    \begin{column}{0.5\textwidth}
      \begin{amberbox}
        \bmheading{bmAmber}{Lo que este proyecto es}
        Un framework de benchmarking reproducible: configuracion declarativa con hash
        canonico, seeds derivadas de una semilla raiz y separacion entre el nucleo del
        benchmark y la logica propia de cada modelo.
      \end{amberbox}
    \end{column}
  \end{columns}
\end{frame}

% ============================================================
% 4. DISENO EXPERIMENTAL
% ============================================================
\begin{frame}{4. Diseno experimental}
  \footnotesize
  \bandedtable
  \begin{tabular}{@{}p{2.1cm}p{6.6cm}p{3.1cm}@{}}
    \toprule
    \textbf{Fase} & \textbf{Pregunta que responde} & \textbf{Estado} \\
    \midrule
    E0A & Cual QCNN de la literatura (Hur, Wei, Gong) tiene mejor desempeno en
          simulacion analitica, sobre los tres datasets del diseno vigente. &
          Completo. 45 corridas. \\
    Quantum Advantage & El QCNN ganador de E0A, supera a una CNN clasica con conteo
          de parametros comparable. & Completo. 30 corridas. \\
    E0B & El QCNN ganador de E0A, supera a una QCNN generica sin sesgo de diseno
          (Cong et al.). & Completo. 30 corridas. \\
    E1 & Como se degrada el desempeno del QCNN ganador al pasar a shots finitos
          (256, 1024, 4096). & Completo. 45 corridas. \\
    E2 & Como se degrada el desempeno bajo ruido de hardware y restricciones de
          conectividad, sobre shots finitos. & Motor construido y validado.
          Matriz completa pendiente. \\
    \bottomrule
  \end{tabular}
\end{frame}

\begin{frame}{4. Diseno experimental: cobertura de modelos por fase}
  \footnotesize
  \bandedtable
  \begin{tabular}{@{}p{2.3cm}p{6.4cm}p{3.1cm}@{}}
    \toprule
    \textbf{Fase} & \textbf{Modelos evaluados} & \textbf{Corridas} \\
    \midrule
    E0A & Los 3 candidatos de la literatura, Hur, Wei y Gong. & 45 \\
    Quantum Advantage & Solo el ganador de E0A, contra su CNN analoga clasica. & 30 \\
    E0B & Solo el ganador de E0A, contra Cong, la QCNN generica de referencia. & 30 \\
    E1 & Solo el ganador de E0A, mas Cong si aplica. & 45 \\
    E2 & Solo el ganador de E0A, mas Cong si aplica. & Pendiente \\
    \bottomrule
  \end{tabular}

  \vspace{8pt}
  \begin{amberbox}
    \bmheading{bmAmber}{Por que las fases posteriores a E0A no repiten los 3 candidatos}
    Con 5 semillas por configuracion, un test de permutacion pareado tiene poca potencia
    para detectar diferencias moderadas. Multiplicar por 3 el numero de comparaciones en
    cada fase posterior diluiria esa potencia sin responder una pregunta nueva. El costo
    de computo tambien crece rapido bajo shots finitos, ver la nota de viabilidad en la
    seccion de proximos pasos.
  \end{amberbox}
\end{frame}

\begin{frame}{4. Diseno experimental: consistencia con el diseno original}
  \footnotesize
  \begin{bluebox}
    \bmheading{bmBlue}{Lo que el diseno original si preveia}
    El regimen de ruido, benchmark-noise en el diseno original, fijaba explicitamente 2
    QCNN, no los 4 modelos de benchmark-core. La practica actual en E2, un ganador mas
    Cong, es consistente con ese punto.
  \end{bluebox}

  \vspace{10pt}

  \begin{amberbox}
    \bmheading{bmAmber}{Lo que cambio despues del diseno original}
    El diseno original, antes de agregar Gong y Cong, preveia correr benchmark-core, E0 y
    E1, con sus 4 modelos de esa version, las 2 QCNN de Hur y Wei mas sus 2 CNN analogas
    clasicas, sin un paso de seleccion de candidato. El embudo de E0A se agrego como
    refinamiento posterior junto con Gong y Cong, y no esta documentado en ese diseno
    original para la fase E1. Se mantiene aqui como decision de diseno confirmada, no
    como continuacion silenciosa de un plan previo.
  \end{amberbox}
\end{frame}

% ============================================================
% 5. ARQUITECTURAS EVALUADAS
% ============================================================
\begin{frame}{\subsec{5.1}{QCNN Hur et al. (2022, Ansatz 8)}}
  \footnotesize
  \begin{columns}[T]
    \begin{column}{0.42\textwidth}
      \begin{bluebox}
        \bmheading{bmBlue}{Especificacion real}
        8 qubits.\\
        36 parametros entrenables.\\
        Encoding: PCA densa a 16 componentes, angulo denso (dos componentes por
        eje de Bloch, ``Angle-compact'').
      \end{bluebox}
    \end{column}
    \begin{column}{0.58\textwidth}
      \begin{indigobox}
        \bmheading{bmIndigo}{Idea central}
        Tres etapas jerarquicas de convolucion y pooling (8 a 4 a 2 a 1 qubits)
        con parametros compartidos entre aplicaciones de una misma etapa,
        replicando la invariancia traslacional de una CNN clasica. Lectura en
        base de Pauli-Z de un solo qubit.
      \end{indigobox}
    \end{column}
  \end{columns}
  \vspace{8pt}
  \begin{amberbox}
    \bmheading{bmAmber}{Por que se incluyo}
    Representa el extremo de codificacion por angulo de la literatura: maximiza
    el numero de features por qubit sin usar codificacion de amplitud.
  \end{amberbox}
\end{frame}

\begin{frame}{\subsec{5.2}{QCNN Wei et al. (2021)}}
  \footnotesize
  \begin{columns}[T]
    \begin{column}{0.42\textwidth}
      \begin{bluebox}
        \bmheading{bmBlue}{Especificacion real}
        10 qubits de trabajo.\\
        46 parametros entrenables (9 del filtro convolucional, 37 del
        hamiltoniano de lectura).\\
        Encoding: amplitud directa (imagen con relleno a $32\times32$, 1024
        amplitudes).
      \end{bluebox}
    \end{column}
    \begin{column}{0.58\textwidth}
      \begin{indigobox}
        \bmheading{bmIndigo}{Idea central}
        Filtro convolucional $3\times3$ como combinacion lineal de unitarios
        (LCU) sobre corrimientos de la imagen. Pooling que descarta 2 de los 10
        qubits. Lectura con un hamiltoniano completo de 37 parametros
        (identidad, Z, ZZ), no un solo observable.
      \end{indigobox}
    \end{column}
  \end{columns}
  \vspace{8pt}
  \begin{amberbox}
    \bmheading{bmAmber}{Por que se incluyo}
    Representa el extremo opuesto a Hur: codificacion de amplitud (maxima
    compresion de informacion por qubit, mayor profundidad de circuito) en vez
    de angulo. Nota de implementacion: el registro de 4 qubits ancilla del
    filtro LCU del paper se reemplaza aqui por su equivalente clasico exacto
    para el caso sin ruido, una combinacion lineal directa, en vez de simular
    puertas controladas y post seleccion real.
  \end{amberbox}
\end{frame}

\begin{frame}{\subsec{5.3}{QCNN Gong et al. (2024)}}
  \footnotesize
  \begin{columns}[T]
    \begin{column}{0.42\textwidth}
      \begin{bluebox}
        \bmheading{bmBlue}{Especificacion real}
        8 qubits.\\
        36 parametros entrenables (mismo Circuit 6 de la familia Hur/Sim et
        al.).\\
        Encoding: tree-hybrid amplitude, bloques de 2 qubits, capacidad de 12
        features (PCA a 8 componentes, relleno con ceros).
      \end{bluebox}
    \end{column}
    \begin{column}{0.58\textwidth}
      \begin{indigobox}
        \bmheading{bmIndigo}{Idea central}
        Combina codificacion de angulo y de amplitud dividiendo los qubits en
        bloques pequenos. Dentro de cada bloque usa rotaciones controladas en
        arbol para aprovechar mejor la profundidad del circuito que el angulo
        denso puro, sin normalizar cada bloque por separado.
      \end{indigobox}
    \end{column}
  \end{columns}
  \vspace{8pt}
  \begin{amberbox}
    \bmheading{bmAmber}{Por que se incluyo}
    Aporta una estrategia de encoding hibrida, intermedia entre Hur (angulo
    denso) y Wei (amplitud pura). Permite evaluar si esa combinacion ofrece
    una ventaja real.
  \end{amberbox}
\end{frame}

\begin{frame}{\subsec{5.4}{QCNN generica de Cong et al. (2019)}}
  \footnotesize
  \begin{columns}[T]
    \begin{column}{0.42\textwidth}
      \begin{bluebox}
        \bmheading{bmBlue}{Especificacion real}
        8 qubits.\\
        51 parametros entrenables (45 de convolucion SU(4), 6 de pooling).\\
        Encoding: angulo simple (PCA a 8 componentes, un feature por qubit).
      \end{bluebox}
    \end{column}
    \begin{column}{0.58\textwidth}
      \begin{indigobox}
        \bmheading{bmIndigo}{Idea central}
        Arquitectura QCNN fundacional y generica. Usa el unitario de
        convolucion SU(4), el mas expresivo posible para 2 qubits, sin ningun
        sesgo de diseno hacia una tarea especifica.
      \end{indigobox}
    \end{column}
  \end{columns}
  \vspace{8pt}
  \begin{amberbox}
    \bmheading{bmAmber}{Por que se incluyo}
    Sirve de referencia en E0B: verifica si la aportacion especifica del
    modelo ganador de E0A (encoding denso, filtro LCU o tree-hybrid) realmente
    aporta ventaja sobre una QCNN jerarquica generica sin ese diseno
    particular.
  \end{amberbox}
\end{frame}

\begin{frame}{\subsec{5.5}{Comparativa de arquitecturas}}
  \footnotesize
  \bandedtable
  \begin{tabular}{@{}lccp{3.4cm}p{3.2cm}@{}}
    \toprule
    \textbf{Modelo} & \textbf{Qubits} & \textbf{Parametros} & \textbf{Encoding} &
      \textbf{Rol en el benchmark} \\
    \midrule
    Hur  & 8  & 36 & PCA-16, angulo denso & Candidato (E0A) \\
    \winner{Wei}  & 10 & 46 & Amplitud directa (1024) & Candidato (E0A), ganador \\
    Gong & 8  & 36 & Tree-hybrid amplitude (PCA-8) & Candidato (E0A) \\
    Cong & 8  & 51 & PCA-8, angulo simple & Referencia generica (E0B) \\
    \bottomrule
  \end{tabular}
\end{frame}

% ============================================================
% 6. RESULTADOS E0A
% ============================================================
\begin{frame}{6. Resultados E0A: media $\pm$ IC 95\% (bootstrap BCa)}
  \footnotesize
  \bandedtable
  \begin{tabular}{@{}lccc@{}}
    \toprule
    \textbf{Dataset} & \textbf{Hur} & \textbf{Wei} & \textbf{Gong} \\
    \midrule
    coat vs shirt & 73.48\% [71.98, 75.10] & \winner{77.10\% [76.68, 77.36]} & 74.10\% [73.42, 74.72] \\
    4 vs 9        & 83.48\% [81.90, 84.18] & \winner{95.38\% [94.76, 95.74]} & 80.92\% [78.60, 83.20] \\
    1 vs 0        & 98.92\% [98.58, 99.18] & \winner{99.24\% [98.96, 99.40]} & 98.74\% [98.64, 98.82] \\
    \bottomrule
  \end{tabular}

  \vspace{10pt}
  Cinco semillas por combinacion. Protocolo identico para los tres modelos: 500
  actualizaciones como techo, lote de 25, early stopping con paciencia de 5
  chequeos.
\end{frame}

\begin{frame}{6. Resultados E0A: honestidad estadistica}
  \begin{amberbox}
    \bmheading{bmAmber}{El piso del test con n=5 semillas}
    La prueba de permutacion pareada de signo, con solo 5 semillas, tiene un
    p-valor minimo posible de 0.0625 antes de corregir por comparaciones
    multiples. Tras la correccion de Holm sobre las 9 comparaciones de E0A,
    ninguna alcanza $p<0.05$, sin importar el tamano del efecto observado.
  \end{amberbox}

  \vspace{8pt}

  \begin{greenbox}
    \bmheading{bmGreen}{Por que se selecciono a Wei de todas formas}
    Wei tiene la media mas alta en los tres datasets. Su intervalo de
    confianza no se traslapa con el de Gong en los tres datasets, y con el de
    Hur en dos de tres (se traslapa en 1 vs 0). El tamano de efecto es grande
    donde mas importa: en 4 vs 9, Wei le gana a Hur y a Gong en 5 de 5
    semillas ($d_z$ entre 5.1 y 6.3, probabilidad de superioridad entre 0.8 y
    1.0). La seleccion se basa en el conjunto de esta evidencia, no en un
    resultado marcado como significativo.
  \end{greenbox}
\end{frame}

% ============================================================
% 7. QUANTUM ADVANTAGE
% ============================================================
\begin{frame}{7. Quantum Advantage: Wei vs. CNN analoga}
  \footnotesize
  \bandedtable
  \begin{tabular}{@{}lccc@{}}
    \toprule
    \textbf{Dataset} & \textbf{Wei (QCNN)} & \textbf{CNN analoga} & \textbf{Traslape de IC} \\
    \midrule
    coat vs shirt & 77.10\% [76.68, 77.36] & 69.40\% [65.56, 72.96] & No \\
    4 vs 9        & 95.38\% [94.76, 95.74] & 79.42\% [69.58, 84.86] & No \\
    1 vs 0        & 99.24\% [98.96, 99.40] & 98.96\% [97.24, 99.62] & Si \\
    \bottomrule
  \end{tabular}

  \vspace{10pt}
  La CNN analoga (\texttt{cnn\_wei\_analoga}) tiene un conteo de parametros
  comparable al de Wei, siguiendo el criterio de Gong et al. Tabla 4.

  \vspace{6pt}
  \begin{amberbox}
    \bmheading{bmAmber}{Significancia}
    Igual que en E0A, con n=5 semillas ninguna comparacion alcanza $p<0.05$
    tras Holm (los p ajustados son 0.19, 0.19 y 0.88). El tamano de efecto en
    4 vs 9 y coat vs shirt es grande ($d_z$ de 1.6 a 1.7, probabilidad de
    superioridad de 1.0), consistente con la ventaja de la QCNN, pero no
    formalmente significativo.
  \end{amberbox}
\end{frame}

% ============================================================
% 8. E0B
% ============================================================
\begin{frame}{8. Validez de la aportacion (E0B): Wei vs. Cong}
  \footnotesize
  \bandedtable
  \begin{tabular}{@{}lccc@{}}
    \toprule
    \textbf{Dataset} & \textbf{Wei} & \textbf{Cong (generica)} & \textbf{Traslape de IC} \\
    \midrule
    coat vs shirt & 77.10\% [76.68, 77.36] & 73.64\% [71.80, 75.12] & No \\
    4 vs 9        & 95.38\% [94.76, 95.74] & 84.02\% [79.87, 87.20] & No \\
    1 vs 0        & 99.24\% [98.96, 99.40] & 98.74\% [98.40, 98.88] & No \\
    \bottomrule
  \end{tabular}

  \vspace{10pt}
  \begin{amberbox}
    \bmheading{bmAmber}{Significancia}
    El intervalo de confianza de Wei no se traslapa con el de Cong en ninguno
    de los tres datasets, y el tamano de efecto es grande en dos de ellos
    ($d_z$ de 1.6 a 2.4). Aun asi, tras la correccion de Holm ninguna
    comparacion alcanza $p<0.05$ (los tres p ajustados son 0.19). La evidencia
    apunta a que la aportacion especifica de Wei si se traduce en ventaja
    sobre una QCNN generica sin ese diseno, pero la confirmacion formal sigue
    limitada por n=5.
  \end{amberbox}
\end{frame}

% ============================================================
% 9. RESULTADOS E1
% ============================================================
\begin{frame}{9. Resultados E1: exactitud de Wei por presupuesto de shots}
  \footnotesize
  \bandedtable
  \begin{tabular}{@{}lcccc@{}}
    \toprule
    \textbf{Dataset} & \textbf{256 shots} & \textbf{1024 shots} & \textbf{4096 shots} &
      \textbf{E0A (referencia)} \\
    \midrule
    coat vs shirt & 63.54\% & 65.34\% & 68.34\% & 77.10\% \\
    4 vs 9        & 80.70\% & 84.26\% & 84.82\% & 95.38\% \\
    1 vs 0        & 98.42\% & 98.62\% & 98.64\% & 99.24\% \\
    \bottomrule
  \end{tabular}

  \vspace{10pt}
  La exactitud sube de forma monotona con el presupuesto de shots en los tres
  datasets. Ninguna comparacion entre presupuestos alcanza $p<0.05$ tras Holm
  (n=5 semillas), aunque el tamano de efecto en coat vs shirt y 4 vs 9 es
  grande ($d_z$ hasta 8.4).

  \vspace{6pt}
  {\scriptsize\color{bmGrey} La columna ``E0A (referencia)'' \textbf{no es
  directamente comparable}: ver la nota metodologica en la siguiente
  diapositiva.}
\end{frame}

\begin{frame}{9. Resultados E1: nota metodologica importante (1/2)}
  \begin{amberbox}
    \bmheading{bmAmber}{Wei perdio capacidad entrenable bajo shots}
    En E0A y en Quantum Advantage, Wei entrena sus 46 parametros originales (9
    del filtro convolucional, 37 del hamiltoniano). Bajo shots finitos, esto
    no fue posible: PennyLane 0.45 no admite diferenciar, con un metodo
    compatible con shots, un \texttt{StatePrep} cuyo vector de entrada depende
    de un parametro entrenable.
  \end{amberbox}

  \vspace{10pt}

  \begin{amberbox}
    \bmheading{bmAmber}{El detalle tecnico}
    La preparacion de estado de Wei (\texttt{MottonenStatePreparation}) exige
    que su vector de entrada tenga norma 1. Perturbar los 9 parametros del
    filtro convolucional para calcular su gradiente rompe esa norma, sin
    importar si el metodo de diferenciacion es parameter-shift o SPSA.
  \end{amberbox}
\end{frame}

\begin{frame}{9. Resultados E1: nota metodologica importante (2/2)}
  \begin{greenbox}
    \bmheading{bmGreen}{La correccion aplicada}
    Los 9 parametros de preparacion de estado se congelan en un valor de
    referencia fijo. En E1 (y en E2) solo se entrenan los 37 parametros del
    hamiltoniano de lectura.
  \end{greenbox}

  \vspace{12pt}

  \begin{indigobox}
    \bmheading{bmIndigo}{Consecuencia para la interpretacion}
    E1 no compara exactamente el mismo modelo que gano E0A, sino una version
    de capacidad reducida (37 de 46 parametros). La degradacion observada
    frente a la columna de referencia de E0A puede deberse, en parte, a esta
    perdida de capacidad y no unicamente al regimen de shots finitos. Separar
    ambos efectos queda pendiente.
  \end{indigobox}
\end{frame}

% ============================================================
% 10. PROXIMOS PASOS
% ============================================================
\begin{frame}{10. Proximos pasos (1/3): E2 y estado de las 4 familias de metricas}
  \begin{bluebox}
    \bmheading{bmBlue}{E2: shots finitos con ruido NISQ}
    El motor de shots, ruido (despolarizante, error de lectura, relajacion
    termica, aisladas y compuestas) y transpilacion a un mapa de acoplamiento
    de 16 qubits ya esta construido y validado con pruebas de plomeria. La
    matriz experimental completa no se ha ejecutado.
  \end{bluebox}

  \vspace{6pt}
  \scriptsize
  \bandedtable
  \begin{tabular}{@{}p{2.2cm}p{6.6cm}p{2.7cm}@{}}
    \toprule
    \textbf{Familia} & \textbf{Estado} & \textbf{Fecha objetivo} \\
    \midrule
    Predictivas & Exactitud desde el inicio. Precision, recall, F1, balanced
      accuracy y macro-F1 ya definidas, aplicadas desde E2. Backfill de E0,
      E0A, Quantum Advantage y E0B, 150 corridas, pendiente. &
      Backfill: fase III \\
    Convergencia & Implementada de forma retroactiva, sin reentrenar. &
      Completo \\
    De recursos & Implementada: qubits, parametros entrenables y presupuesto
      de shots por modelo. & Completo \\
    Robustez & Bloqueada, depende de la matriz de E2. & Con E2 \\
    \bottomrule
  \end{tabular}
\end{frame}

\begin{frame}{10. Proximos pasos (2/3): viabilidad de expandir E1 a Hur y Gong}
  \footnotesize
  Se cronometraron actualizaciones de gradiente reales, parameter-shift, no
  backprop, de Hur y Gong bajo 256 shots, el presupuesto mas bajo de E1,
  sobre datos reales, antes de comprometerse a una matriz completa.

  \vspace{6pt}
  \bandedtable
  \begin{tabular}{@{}lcc@{}}
    \toprule
    \textbf{Modelo} & \textbf{Segundos por actualizacion} &
      \textbf{Matriz completa estimada (45 corridas)} \\
    \midrule
    Hur  & 62.2 s & 389 h, unos 16 dias \\
    Gong & 21.0 s & 131 h, unos 5 dias y medio \\
    \bottomrule
  \end{tabular}

  \vspace{6pt}
  \begin{amberbox}
    \bmheading{bmAmber}{Por que Wei, ya ejecutado, es tan distinto}
    Wei diferencia coeficientes de un hamiltoniano de lectura, una expansion
    lineal barata. Hur y Gong diferencian los parametros del circuito
    completo via parameter-shift, ordenes de magnitud mas caro bajo shots.
  \end{amberbox}

  \vspace{6pt}
  \textbf{Recomendacion}: con el protocolo actual, 500 actualizaciones, lote
  de 25, una expansion v1.1 de E1 a los 3 candidatos no es viable dentro de
  la fase III. Es viable solo si se reduce el numero de actualizaciones, el
  tamano de lote, o el barrido de presupuestos de shots, o si se paraleliza
  en mas de una maquina.
\end{frame}

\begin{frame}{10. Proximos pasos (3/3): calendario planeado vs. ejecutado}
  \scriptsize
  \bandedtable
  \begin{tabular}{@{}p{3.6cm}p{2.4cm}p{6.4cm}@{}}
    \toprule
    \textbf{Hito (plan de trabajo)} & \textbf{Fecha planeada} & \textbf{Estado real} \\
    \midrule
    Benchmark E0 completo & 2 sep 2026 & Completado, 27 ago 2026, solo exactitud \\
    Backfill de metricas predictivas de E0 & No estaba en el plan original & En curso desde 21 sep 2026, ETA unas 6 h \\
    Benchmark E0-E1 completo & 9 sep 2026 & E1 (solo Wei) completado 17 sep 2026. Ampliado a Hur+Wei+Gong: en curso desde 21 sep 2026, ETA no antes del 13 oct 2026 (~22 dias, ver diapositiva anterior) \\
    Motor NISQ (E2) validado & 17 sep 2026 & Validado. Bug de transpilacion para Cong encontrado y corregido 18 sep 2026 \\
    Benchmark de ruido ejecutado & 23 sep 2026 & No inicia hasta terminar la expansion de E1 (orden confirmado 21 sep 2026). ETA no antes de mediados de oct 2026 \\
    \bottomrule
  \end{tabular}

  \vspace{8pt}
  Las fechas de E0 y E1 (primera corrida) provienen de las marcas de tiempo
  de sus archivos de provenance. Las fechas ETA son proyecciones a partir
  de costo medido, no garantias: el early stopping puede acortarlas algo,
  no cambia el orden de magnitud. No se propone aqui una fecha de
  recuperacion para el 23 de sep: se documenta el desfase y su causa.
\end{frame}

% ============================================================
% 11. CIERRE
% ============================================================
\begin{frame}{11. Cierre}
  \vspace{6pt}
  \begin{indigobox}
    \bmheading{bmIndigo}{Agradecimientos}
    Al Dr. Leonardo Ledesma Dominguez, director de este proyecto, y al
    programa Next Generation Scientists (NGS).
  \end{indigobox}

  \vspace{12pt}

  \bmheading{bmIndigo}{Referencias principales}
  \begin{itemize}
    \small
    \item Hur, T., Kim, L., Park, D. K. (2022). Quantum convolutional neural
      network for classical data classification. \textit{Quantum Machine
      Intelligence}.
    \item Wei, S., Chen, Y., Zhou, Z., Long, G. (2021). A quantum
      convolutional neural network on NISQ devices. \textit{arXiv:2104.06918}.
    \item Gong, L.-H., Pei, J.-J., Zhang, T.-F., Zhou, N.-R. (2024). Quantum
      convolutional neural network based on variational quantum circuits.
      \textit{Optics Communications}, 550, 129993.
    \item Cong, I., Choi, S., Lukin, M. D. (2019). Quantum convolutional
      neural networks. \textit{Nature Physics}, 15, 1273-1278.
  \end{itemize}
\end{frame}

\end{document}
